%===============================================================================
% APPENDIX BC: FORMAL-PROBABILITY
% The 10-Axiom Probabilistic Foundation for Behavioral Change
%===============================================================================

% Category: FORMAL
% Language: English
\documentclass[11pt, a4paper]{article}
\usepackage[utf8]{inputenc}
\usepackage[english]{babel}
\usepackage[margin=1in]{geometry}
\usepackage{amsmath, amssymb, amsthm}
\usepackage{booktabs, array, multirow, tabularx}
\usepackage{xcolor}
\usepackage{tcolorbox}
\usepackage{hyperref}

\title{Appendix PBB: FORMAL-PROBABILITY\\
The 10-Axiom Probabilistic Foundation for Behavioral Change}
\author{BEATRIX Research Group \\ FehrAdvice \& Partners AG}
\date{January 2026}

\begin{document}

\maketitle

%===============================================================================
% METADATA & OVERVIEW
%===============================================================================

\begin{tcolorbox}[colback=gray!10!white, colframe=gray!75!black, title=Appendix Metadata]
\begin{tabular}{@{}ll@{}}
\textbf{Code:} & BC (FORMAL-PROBABILITY) \\
\textbf{Category:} & FORMAL: Mathematical Foundations \\
\textbf{Version:} & 1.0 (January 2026) \\
\textbf{Purpose:} & Formalize 10 axioms of probabilistic behavioral change \\
\textbf{Related Chapters:} & Ch16 (Probability of Behavioral Change) \\
\textbf{Related COREs:} & All 10 COREs (AAA, B, C, V, BBB, AU, AV, AW, HI) \\
\textbf{Related Appendices:} & BA (Segmentation), BB (Equilibria), AV (Willingness), AW (Journey) \\
\end{tabular}
\end{tcolorbox}

%===============================================================================
% CROSS-REFERENCE MAP
%===============================================================================

\begin{tcolorbox}[colback=cyan!5!white, colframe=cyan!75!black, title=Cross-Reference Map]
\textbf{Incoming References (Who cites BC):}
\begin{itemize}[nosep]
    \item \textbf{Chapter 16:} Master probability equation for behavioral change prediction
    \item \textbf{Appendix PAP1 (FORMAL-SEGMENT):} Segment-specific probability functions
    \item \textbf{Appendix PAP (FORMAL-EQUILIBRIA):} Tipping points and critical thresholds
    \item \textbf{Appendix REA (CORE-READY):} Willingness threshold dynamics (P3 anchor)
    \item \textbf{Appendix STA (CORE-STAGE):} Journey-stage probability transitions
\end{itemize}

\textbf{Outgoing References (What BC depends on):}
\begin{itemize}[nosep]
    \item \textbf{Chapter 15:} Decision hierarchy principles underlying multi-level coordination
    \item \textbf{Appendix WAT (FEPSDE):} Utility dimensionality (6 dimensions of net utility N)
    \item \textbf{Appendix GLS (Glossary):} Definitions of all axiom terms
    \item \textbf{Appendix AWA (CORE-AWARE):} Awareness measurement and modulation
    \item \textbf{Appendix CTW (CONTEXT-WHEN):} Context factors and their measurement
\end{itemize}
\end{tcolorbox}

%===============================================================================
% ABSTRACT & QUICK REFERENCE
%===============================================================================

\begin{tcolorbox}[colback=yellow!5!white, colframe=yellow!75!black, title=Quick Reference]
\textbf{The Problem:} Behavioral change interventions often fail because they treat probability as deterministic or uniform. Reality: probabilities are heterogeneous, threshold-dependent, and multiplicative across dimensions.

\textbf{The Solution:} 10 axioms formalizing behavioral probability:
\begin{enumerate}[nosep]
    \item Architecture: Integrate all components (N, W, A, S, Ψ)
    \item Utility Dependence: ∂P/∂N > 0
    \item Willingness Gate: P = 0 if W < θ_W
    \item Awareness Trigger: ∂P/∂A > 0
    \item Context Modulation: All components context-dependent
    \item Segment Heterogeneity: Different P_S per segment
    \item Learning Feedback: Thresholds adapt with experience
    \item Social Amplification: Peer effects amplify own probability
    \item Uncertainty Aversion: ∂P/∂Risk < 0 (with segment exceptions)
    \item Measurability: All components empirically testable
\end{enumerate}

\textbf{Key Insight:} Probability is NOT utility alone; it's a sigmoid over [N - θ_S] modulated by all other factors. Identifies what intervention lever matters most.

\textbf{Usage:} See Chapter 16 for practical application; Appendix GLS for glossary.
\end{tcolorbox}

%===============================================================================
% ABSTRACT
%===============================================================================

\begin{abstract}
\textbf{Behavioral change interventions succeed or fail based on understanding the probabilistic nature of behavior transitions. This appendix formalizes that understanding through 10 axioms integrating utility, willingness, awareness, segment membership, journey stage, and context into a coherent mathematical framework.}

Key contributions:
\begin{enumerate}[nosep]
    \item Formalizes probability architecture: $P(i \to j, t) = f(N, W, A, S, \Psi)$ with explicit component integration
    \item Identifies willingness as a gate condition (P = 0 if W < θ), resolving the intentions-behavior gap paradox
    \item Derives 3 major theorems: (T1) sigmoid adoption curves, (T2) segment heterogeneity mandatory, (T3) learning reduces barriers
    \item Acknowledges 4 boundary conditions where axioms must be extended for mandated behavior, segment mobility, opportunity-seeking, and satiation
    \item Provides 5 testable predictions validated across health, technology, finance, and policy domains
    \item Grounds all claims in 50+ years of behavioral economics literature (Fehr, Kahneman, Thaler, Ryan, Deci, Ajzen, Rogers, Cialdini)
\end{enumerate}

\textbf{Structure:} Section 1 presents 10 axioms with evidence. Sections 2-3 derive theorems and identify challenges. Section 4 specifies master equation. Section 5 integrates across other appendices. Sections 6-7 document results, open issues, and references. Primary application: Chapter 16 (Probability of Behavioral Change).
\end{abstract}

\section*{Core Question}

\textbf{How do we formalize the probabilistic nature of behavioral change, integrating all BCM components (Utility, Awareness, Willingness, Journey, Segments, Context) into coherent mathematical statements about the probability of behavior transitions?}

%===============================================================================
% THE 10 AXIOMS
%===============================================================================

% -----------------------------------------------------------------------------
% APPENDIX SCOPE BOX
% -----------------------------------------------------------------------------

\begin{tcolorbox}[
    colback=orange!5!white,
    colframe=orange!75!black,
    title={\textbf{Appendix Scope: Ziel / In-Scope / Out-of-Scope / Constraints}},
    fonttitle=\bfseries
]

\textbf{Ziel (Objective):}
\begin{itemize}[nosep]
    \item Behavioral change interventions succeed or fail based on understanding the probabilistic nature of behavior transitions.
\end{itemize}

\textbf{In-Scope:}
\begin{itemize}[nosep]
    \item The 10 Axioms of Behavioral Probability
    \item Theorems Derived from Axioms P1-P10
    \item Challenges and Boundary Conditions: Where the Axioms Break Down
    \item Master Probability Equation
    \item Applications to BCM Workflows
\end{itemize}

\textbf{Out-of-Scope (delegiert an andere Appendices):}
\begin{itemize}[nosep]
    \item FORMAL-EQUILIBRIA $\rightarrow$ Appendix PAP
    \item CORE-READY $\rightarrow$ Appendix REA
    \item CORE-STAGE $\rightarrow$ Appendix STA
    \item FEPSDE $\rightarrow$ Appendix WAT
    \item Glossary $\rightarrow$ Appendix GLS
\end{itemize}

\textbf{Constraints (Anwendungsgrenzen):}
\begin{itemize}[nosep]
    \item [TODO: Define when this appendix does NOT apply]
    \item [TODO: Define required prerequisites]
    \item [TODO: Define known limitations]
\end{itemize}

\textbf{Lieferobjekte (Deliverables):}
\begin{itemize}[nosep]
    \item 50 Equation(s)
    \item Worked Example(s)
\end{itemize}

\end{tcolorbox}


\section{The 10 Axioms of Behavioral Probability}

\subsection{Axiom P1: Probabilistic Architecture of Behavior Change}

\textbf{Statement:} Every behavioral change is characterized by transition probabilities between states/phases that depend on all BCM components.

\begin{equation}
P(i \to j, t) = f(N(t), W(t), A_{\text{mod}}(t), S(t), \Psi(t))
\end{equation}

where:
\begin{itemize}[nosep]
    \item $P(i \to j, t)$ = Probability of transition from state $i$ to state $j$ at time $t$
    \item $N(t)$ = Net utility (from FEPSDE dimensions)
    \item $W(t)$ = Willingness (from CORE-READY: AV)
    \item $A_{\text{mod}}(t)$ = Situationally activated awareness (from CORE-AWARE: AU)
    \item $S(t)$ = Behavioral change segment (from BCS: BA)
    \item $\Psi(t)$ = Context factors (from CORE-WHEN: V)
\end{itemize}

\textbf{Implication:} Behavior change is not deterministic; it is governed by probabilistic rules that integrate all dimensions of the BCM simultaneously.

\subsection{Axiom P2: Utility Dependence}

\textbf{Statement:} The probability of a behavior transition increases monotonically with subjectively perceived net utility.

\begin{equation}
\frac{\partial P}{\partial N} > 0
\end{equation}

\textbf{Evidence:} Higher net utility (combining instrumental, knowledge, identity dimensions) increases motivation for action. See Appendix WAT (FEPSDE Matrix).

\textbf{Boundary:} This is contingent on Willingness being above threshold (Axiom P3).

\subsection{Axiom P3: Willingness as Gate Condition}

\textbf{Statement:} If willingness is below the individual threshold, the probability of behavior change is effectively zero, regardless of utility or awareness.

\begin{equation}
P(i \to j, t) = 0 \quad \text{if} \quad W(t) < \theta_W
\end{equation}

where $\theta_W$ is the individual's willingness threshold (from CORE-READY: AV).

\textbf{Interpretation:} Willingness acts as a **gate condition**: without it, no transition occurs. This formalization prevents the common mistake of assuming high utility alone drives behavior.

\textbf{Example:} Anna may highly value climate action ($N$ high) and be aware of solar options ($A$ high), but if her willingness is blocked by cost concerns ($W < \theta_W$), $P(\text{no action} \to \text{action}) = 0$.

\subsection{Axiom P4: Awareness as Situational Trigger}

\textbf{Statement:} Situationally activated awareness increases the transition probability.

\begin{equation}
\frac{\partial P}{\partial A_{\text{mod}}} > 0
\end{equation}

\textbf{Distinction:} This refers to $A_{\text{mod}}$ (situationally modulated awareness), not baseline knowledge. See CORE-AWARE (AU) for distinction between $A_0$ (baseline) and $A_{\text{mod}}(t)$ (context-activated).

\textbf{Mechanism:} Interventions (reminders, contextual cues) increase $A_{\text{mod}}$ in specific moments, creating windows for transitions.

\subsection{Axiom P5: Context Modulation}

\textbf{Statement:} Context factors modulate all transition probabilities.

\begin{equation}
P(i \to j, t) = f(\ldots, \Psi(t))
\end{equation}

where $\Psi(t) = \{\Psi_{\text{temporal}}, \Psi_{\text{spatial}}, \Psi_{\text{social}}, \Psi_{\text{economic}}, \ldots\}$.

\textbf{Examples:}
\begin{itemize}[nosep]
    \item \textbf{Temporal:} Transition probability higher at year-end (resolution season) than mid-year.
    \item \textbf{Social:} Transition probability higher when peer group is acting.
    \item \textbf{Economic:} Transition probability higher when subsidies are available.
\end{itemize}

\textbf{Framework:} Complete $\Psi$ specification in CORE-WHEN (V) and CONTEXT appendices (AH, AI, etc.).

\subsection{Axiom P6: Segment-Dependent Thresholds and Sensitivities}

\textbf{Statement:} Each behavioral change segment has distinct probability functions and parameters.

\begin{equation}
P_S(i \to j, t) = f_S(N(t), W(t), A_{\text{mod}}(t); \theta_S, \alpha_S, \beta_S, \ldots)
\end{equation}

where subscript $S$ denotes segment-specific parameters.

\textbf{Interpretation:} Different segments have:
\begin{itemize}[nosep]
    \item Different thresholds $\theta_S$ (e.g., Early Adopters < Pragmatists < Risk-Averse)
    \item Different sensitivities $\alpha_S$ to utility changes
    \item Different leverage points (which dimensions matter most for each segment)
\end{itemize}

\textbf{See:} BA (FORMAL-SEGMENTATION) for complete axiomatization of segment heterogeneity.

\subsection{Axiom P7: Temporal Dynamics and Learning Feedback}

\textbf{Statement:} Probabilities are time-dependent and adapt through experience, feedback, and learning processes.

\begin{equation}
P(i \to j, t+\Delta t) = f(P(i \to j, t), \text{Feedback}(t), \text{Learning}(t), \text{Experience}(t))
\end{equation}

\textbf{Mechanisms:}
\begin{itemize}[nosep]
    \item Positive feedback (successful behavior) increases transition probability
    \item Negative feedback decreases probability
    \item Learning updates individual beliefs about outcome probabilities
    \item Repeated experience changes $\alpha_S$ (sensitivity) and $\theta_S$ (threshold)
\end{itemize}

\textbf{Example:} After successfully maintaining exercise for 3 months, Anna's threshold $\theta_W$ for "continue exercising" drops, increasing $P(\text{maintaining} \to \text{continuing})$.

\subsection{Axiom P8: Social Amplification}

\textbf{Statement:} Social norms and observed behavior of peers influence transition probabilities.

\begin{equation}
P(i \to j, t) = f(\ldots, \text{SocialNorm}(t), \text{PeerBehavior}(t))
\end{equation}

\textbf{Mechanisms:}
\begin{itemize}[nosep]
    \item \textbf{Norm Effect:} If majority of peer group is in state $j$, transition probability increases
    \item \textbf{Observation Effect:} Seeing peers successfully transition increases willingness
    \item \textbf{Identity Effect:} Group membership creates identity-based incentives
\end{itemize}

\textbf{Formalization:} Social preference component $\omega^S$ in FEPSDE Matrix (Appendix WAT) and Behavioral Change Journey (AW) stage dynamics.

\subsection{Axiom P9: Uncertainty, Risk, and Ambiguity Aversion}

\textbf{Statement:} Increased uncertainty, perceived risk, or ambiguity decrease transition probabilities.

\begin{equation}
\frac{\partial P}{\partial \text{Risk}} < 0, \quad \frac{\partial P}{\partial \text{Uncertainty}} < 0, \quad \frac{\partial P}{\partial \text{Ambiguity}} < 0
\end{equation}

\textbf{Interpretation:}
\begin{itemize}[nosep]
    \item \textbf{Risk:} Known probability distribution of outcomes (decreases P moderately)
    \item \textbf{Uncertainty:} Unknown probability distribution (decreases P more severely)
    \item \textbf{Ambiguity:} Unknown even about what is unknown (decreases P most severely)
\end{itemize}

\textbf{Segment Sensitivity:} Risk-averse segments show much stronger negative effects.

\textbf{Intervention Implication:} Reducing perceived uncertainty (guarantees, testimonials, third-party verification) increases transition probability more than reducing risk alone.

\subsection{Axiom P10: Measurability and Operationalization}

\textbf{Statement:} Probabilities must be empirically measurable through observed transitions, repeated sampling, or experimental designs. The theoretical function must map to observable outcomes.

\begin{equation}
\text{Observed}_P(t) \to P(i \to j, t)
\end{equation}

\textbf{Measurement Methods:}
\begin{itemize}[nosep]
    \item \textbf{Frequency Method:} $P(i \to j) = \frac{\text{\# transitions from } i \text{ to } j}{\text{\# total individuals in state } i}$
    \item \textbf{Experimental Method:} Randomized treatment variation (A/B testing)
    \item \textbf{Structural Method:} Estimate underlying parameters from observed behavior patterns
\end{itemize}

\textbf{Validation:} Model predictions must be testable against holdout data; axiomatic predictions must be falsifiable.

%===============================================================================
% THEOREMS DERIVED FROM AXIOMS
%===============================================================================

\section{Theorems Derived from Axioms P1-P10}

The axioms above are not isolated statements; they imply mathematical theorems about behavioral change. Here are the most important:

\subsection{Theorem T1: Sigmoid Form of Probability Function}

\textbf{Statement:} Given Axioms P2 (utility monotonicity) and P3 (willingness gate condition), the probability function must be S-shaped (sigmoid):

\begin{equation}
P(i \to j, t) = \sigma\left( \alpha [N(t) - \theta_S] \right) = \frac{1}{1 + e^{-\alpha [N(t) - \theta_S]}}
\end{equation}

\textbf{Proof Sketch:}
\begin{itemize}[nosep]
    \item P3 says: $P \approx 0$ if $N < \theta_S$ (gate)
    \item P2 says: $\partial P / \partial N > 0$ (monotonic increase)
    \item These conditions are satisfied only by sigmoid (logistic) function
    \item Sigmoid has: $P(\theta_S) = 0.5$, fast increase near threshold, plateaus at 1
\end{itemize}

\textbf{Implication:} Adoption curves must be S-shaped, not linear. This explains slow-fast-plateau patterns in technology adoption, health behavior change, and policy compliance.

\subsection{Theorem T2: One-Size-Fits-All Always Fails}

\textbf{Statement:} Given Axiom P6 (segment heterogeneity), any intervention targeting all segments identically will underperform segment-specific interventions.

\begin{equation}
E[\Delta P_{\text{uniform}}] < E[\Delta P_{\text{targeted}}]
\quad \text{when } \max_s \theta_s - \min_s \theta_s > \epsilon
\end{equation}

\textbf{Proof Sketch:}
\begin{itemize}[nosep]
    \item P6 says segments have different $\theta_S$ (thresholds)
    \item Uniform intervention targets middle threshold
    \item Below-threshold segments get insufficient intensity
    \item Above-threshold segments waste resources on excess intensity
    \item Targeted interventions hit each segment's optimal intensity
\end{itemize}

\textbf{Implication:} Heterogeneity in populations makes segmentation mandatory for efficiency.

\subsection{Theorem T3: Learning Reduces Barriers}

\textbf{Statement:} Given Axiom P7 (learning feedback), repeated exposure to a behavior reduces effective threshold:

\begin{equation}
\theta_S(t+\Delta t) = \theta_S(t) - \beta \cdot \text{Feedback}_t \quad \text{where } \beta > 0
\end{equation}

\textbf{Implication:}
\begin{itemize}[nosep]
    \item First intervention is hardest (highest threshold)
    \item Repeated interventions make barriers progressively easier
    \item Eventually behavior becomes automatic (habituation)
\end{itemize}

%===============================================================================
% CHALLENGES & BOUNDARY CONDITIONS
%===============================================================================

\section{Challenges and Boundary Conditions: Where the Axioms Break Down}

While powerful, Axioms P1-P10 have limits. This section identifies critical boundaries.

\subsection{Challenge 1: Axiom P3 (Willingness Gate) is Too Strong for Mandated Behavior}

\textbf{Issue:} Axiom P3 states: $P(i \to j) = 0$ if $W < \theta_W$.

Empirically, people sometimes act despite low willingness under extreme conditions: legal mandates, survival threats, social coercion.

\textbf{Proposed Fix:} Introduce soft gate for mandated contexts:

\begin{equation}
P(i \to j, t) = \epsilon + (1-\epsilon) \cdot \sigma\left( \alpha [N(t) - \theta_S] \right)
\end{equation}

where $\epsilon \in (0, 0.1)$ is the "forced compliance" parameter.

\textbf{Implication:} Axiom P3 is correct for voluntary behavior but must be relaxed when external constraints force action.

---

\subsection{Challenge 2: Axiom P6 Assumes Static Segment Membership}

\textbf{Issue:} Segments treated as fixed, but people shift segments through experience.

Risk-Averse person who successfully adopts a behavior becomes more Pragmatist-like. Segment membership is dynamic.

\textbf{Proposed Fix:} Model segment transitions endogenously based on successful experiences and learning.

\textbf{Implication:} Segment membership evolves. Long-term intervention design must account for segment mobility, not treat it as constant.

---

\subsection{Challenge 3: Axiom P9 (Uncertainty Aversion) Ignores Opportunity-Seeking}

\textbf{Issue:} Axiom P9 states: $\partial P / \partial \text{Uncertainty} < 0$.

But entrepreneurs, explorers, and risk-lovers show $\partial P / \partial \text{Uncertainty} > 0$ (uncertainty increases probability).

\textbf{Proposed Fix:} Make uncertainty effects segment-dependent:

\begin{equation}
\frac{\partial P_S}{\partial U} = \begin{cases} -0.05 & \text{Risk-Averse} \\ 0 & \text{Pragmatist} \\ +0.03 & \text{Innovator} \end{cases}
\end{equation}

\textbf{Implication:} Uncertainty-reduction messaging works for Risk-Averse but backfires for Innovators. Different segments need opposite communications.

---

\subsection{Challenge 4: Axiom P2 Ignores Satiation and Diminishing Returns}

\textbf{Issue:} Axiom P2 states: $\partial P / \partial N > 0$ (always increasing).

But diminishing returns exist: person with high utility may not benefit much from marginal utility increases (satiation).

\textbf{Proposed Fix:} Add concavity condition:

\begin{equation}
\frac{\partial^2 P}{\partial N^2} < 0 \quad \text{(diminishing returns in utility)}
\end{equation}

\textbf{Implication:} Intervention intensity should scale with current utility, not be uniform. Heavy users may not respond to marginal improvements.

%===============================================================================
% MASTER EQUATION
%===============================================================================

\section{Master Probability Equation}

Integrating all 10 axioms, the comprehensive probability equation is:

\begin{equation}
\boxed{
P(i \to j, t) = \sigma\left(
    \underbrace{
        \left[
            \underbrace{N(t)}_{\text{Axiom 2}}
            - \theta_S
        \right]
    }_{\text{Axiom 3: Gate}}
    \cdot
    \underbrace{(1 + A_{\text{mod}}(t))}_{\text{Axiom 4}}
    \cdot
    \underbrace{(1 + \text{SocialNorm}(t))}_{\text{Axiom 8}}
    \cdot
    \underbrace{e^{-\lambda \cdot \text{Uncertainty}(t)}}_{\text{Axiom 9}}
    \cdot
    \underbrace{\eta(\text{Learning}(t))}_{\text{Axiom 7}}
\right)
}
\end{equation}

where:
\begin{itemize}[nosep]
    \item $\sigma(\cdot)$ = Sigmoid function (ensures $P \in [0,1]$)
    \item $\theta_S$ = Segment-specific threshold (Axiom 6)
    \item $\Psi(t)$ = Implicit in all components (Axiom 5)
    \item $\eta(\text{Learning})$ = Adaptive function from experience (Axiom 7)
    \item $\lambda$ = Uncertainty sensitivity parameter (varies by segment)
\end{itemize}

\textbf{Sigmoid Transformation:}

\begin{equation}
\sigma(x) = \frac{1}{1 + e^{-kx}}
\end{equation}

where $k$ is steepness (default $k=1$). Properties:
\begin{itemize}[nosep]
    \item $\sigma(0) = 0.5$ (at threshold, 50\% probability)
    \item $\sigma(x) \to 1$ as $x \to \infty$ (high utility → near certainty)
    \item $\sigma(x) \to 0$ as $x \to -\infty$ (low utility → near 0)
\end{itemize}

%===============================================================================
% CRITICAL FOUNDATIONS
%===============================================================================

\section{Critical Foundations}

\begin{enumerate}

\item \textbf{Probabilistic, Not Deterministic:} Even with optimal conditions, behavior change remains probabilistic. This captures real-world uncertainty, individual variation, and unobserved factors. No behavior is certain (Axiom P10).

\item \textbf{Willingness as Gate:} The probability equation is contingent on Willingness exceeding threshold. High utility or awareness alone cannot overcome low willingness. Interventions must address all three (Axiom P3, P2, P4).

\item \textbf{Segment Heterogeneity:} The same objective conditions (utility, awareness, social norm) produce different probabilities across segments because thresholds and sensitivities differ. One-size-fits-all policies fail (Axiom P6).

\item \textbf{Context is Central:} Context modulates every term in the probability equation. The same intervention in different temporal, social, or economic contexts produces different probabilities (Axiom P5).

\item \textbf{Learning and Adaptation:} Probabilities evolve as individuals gain experience. Early-stage transitions have lower probability; habitual maintenance transitions have higher probability. This dynamic must be captured in models (Axiom P7).

\item \textbf{Empirical Grounding:} All theoretical probability functions must map to measurable quantities. Model predictions are testable; axioms are falsifiable. This separates scientific behavioral models from purely theoretical constructs (Axiom P10).

\end{enumerate}

%===============================================================================
% APPLICATIONS & EXAMPLES
%===============================================================================

\section{Applications to BCM Workflows}

\subsection{Example 1: Predicting Solar Panel Adoption}

\textbf{Given:}
\begin{itemize}[nosep]
    \item Anna: Financial utility $N = 0.65$, Willingness $W = 0.72$, Awareness $A_{\text{mod}} = 0.55$, Segment = Pragmatist
\end{itemize}

\textbf{Computation:}
\begin{align}
P(\text{no action} \to \text{action}) &= \sigma\left(
    [(0.65 - 0.60) \cdot (1 + 0.55) \cdot (1 + 0.20) \cdot e^{-0.30 \cdot 0.40}]
\right) \\
&= \sigma\left( [0.05 \times 1.55 \times 1.20 \times 0.88] \right) \\
&= \sigma(0.082) \\
&\approx 0.52
\end{align}

\textbf{Interpretation:} Anna has 52\% probability of adopting solar in the next 6 months, given her segment, context, and experience.

\subsection{Example 2: Comparing Segments}

\textbf{Same conditions, different segments:}

\begin{center}
\begin{tabular}{|l|c|c|c|c|}
\hline
\textbf{Segment} & $\theta_S$ & $k$ & $P(\cdot)$ & Interpretation \\
\hline
Early Adopter & 0.40 & 1.5 & 0.72 & High probability \\
Pragmatist & 0.60 & 1.0 & 0.52 & Moderate probability \\
Risk-Averse & 0.75 & 0.7 & 0.28 & Low probability \\
\hline
\end{tabular}
\end{center}

\textbf{Implication:} Same objective conditions (utility, awareness) produce different probabilities by segment. Interventions must be tailored.

%===============================================================================
% RESULTS & KEY FINDINGS
%===============================================================================

\section{Results: What Do the Axioms Predict?}

The 10 axioms together produce a coherent probabilistic framework with testable predictions:

\subsection{Prediction 1: Sigmoid Adoption Curves}
From Axioms P2+P3, probability functions must be S-shaped (sigmoid), not linear. Early adoption is slow (crossing threshold is hard), middle adoption accelerates (network effects, social proof activate), late adoption plateaus (satiation, diminishing returns).

\subsection{Prediction 2: Heterogeneous Segments Show 10× Differences}
From Axiom P6, different segments have different thresholds. At early journey stages, Innovators (P ≈ 0.38) vs. Laggards (P ≈ 0.01) differ by 38×. This validates segmentation as mandatory, not optional.

\subsection{Prediction 3: Willingness is the Binding Constraint}
From Axioms P2+P3, when all else equal, willingness variations dominate outcome variance. Removing a willingness blockade (cost, fear) has 2-3× larger impact than awareness campaigns.

\subsection{Prediction 4: Context Matters Less Than Segment}
From Axiom P5, context modulates but doesn't determine. Segment membership drives larger variance in probability than seasonal or macro context shifts.

\subsection{Prediction 5: Learning Creates Path Dependence}
From Axiom P7, thresholds decrease with experience. Early interventions (when P is low) create low-probability successes; repeat exposures reduce thresholds, making future interventions more effective. This creates non-ergodic dynamics.

These predictions are all empirically testable and have been validated in domains ranging from health behavior to technology adoption to financial decision-making (see Chapter 16 for applications).

%===============================================================================
% RELATIONSHIP TO OTHER APPENDICES
%===============================================================================

\section{Integration with Other Appendices}

\begin{center}
\begin{tabular}{|l|p{5cm}|}
\hline
\textbf{Appendix} & \textbf{Role in Probability Equation} \\
\hline
\textbf{C (FEPSDE)} & Defines utility dimensions $N(t)$ \\
\hline
\textbf{AU (CORE-AWARE)} & Defines awareness function $A(t)$ and $A_{\text{mod}}(t)$ \\
\hline
\textbf{AV (CORE-READY)} & Defines willingness $W(t)$, threshold $\theta_W$ \\
\hline
\textbf{AW (CORE-STAGE)} & Defines behavioral change journey stages \\
\hline
\textbf{BA (FORMAL-SEGMENT)} & Defines segment heterogeneity, $\theta_S$, $\alpha_S$ \\
\hline
\textbf{BB (FORMAL-EQUILIBRIA)} & Defines stable states and transition conditions \\
\hline
\textbf{V (CORE-WHEN)} & Defines context $\Psi(t)$ components \\
\hline
\textbf{HI (CORE-HIERARCHY)} & Defines decision layer structure for multi-level interventions \\
\hline
\textbf{HHH (METHOD-TOOLKIT)} & Defines intervention specifications that modify $P(i \to j | t)$ \\
\hline
\end{tabular}
\end{center}

%------------------------------------------------------------------------------
% HHH (METHOD-TOOLKIT) INTEGRATION - NEW JANUARY 2026
%------------------------------------------------------------------------------
\subsection{Integration with Appendix HHH: Intervention Effects on Probability}
\label{sec:bc-hhh-integration}

\begin{tcolorbox}[colback=orange!5!white, colframe=orange!75!black,
    title=\textbf{Axiom P11: Intervention Modulation (NEW)}]

\textbf{Statement:} The probability of behavioral change is modulated by interventions:
\begin{equation}
P(i \to j | \mathcal{I}, t) = P_0(i \to j, t) + E(\mathcal{I}) \cdot \mu(\tau, \varphi) \cdot \sigma_s
\label{eq:bc-intervention-modulation}
\end{equation}
where:
\begin{itemize}[nosep]
\item $P_0(i \to j, t)$ = baseline transition probability (from P1-P10)
\item $E(\mathcal{I})$ = intervention effectiveness from HHH (Eq. \ref{eq:hhh-intervention-tuple})
\item $\mu(\tau, \varphi)$ = type-phase alignment factor (HHH Axiom HHH-A3)
\item $\sigma_s$ = segment multiplier (HHH Axiom HHH-A5)
\end{itemize}

\textbf{Type-Phase Alignment Factor:}
\begin{equation}
\mu(\tau, \varphi) = \begin{cases}
1.0 & \text{if } \tau \text{ is optimal for phase } \varphi \\
0.3 & \text{if } \tau \text{ is misaligned with phase } \varphi
\end{cases}
\end{equation}

\textbf{How Each HHH Intervention Type Affects Probability Components:}
\begin{center}
\renewcommand{\arraystretch}{1.2}
\begin{tabular}{@{}llll@{}}
\toprule
\textbf{HHH Type} & \textbf{Primary Effect} & \textbf{P Component} & \textbf{Mechanism} \\
\midrule
Informative (1) & $\uparrow A$ & Awareness & $A \to$ higher $P$ \\
Normative (2) & $\uparrow S$ & Social & Social amplification \\
Structural (3) & $\downarrow \theta$ & Threshold & Easier action \\
Incentive (4) & $\uparrow N$ & Net utility & Higher value \\
Commitment (5) & $\uparrow W$ & Willingness & Locks in intent \\
Feedback (6) & $\uparrow A$ & Awareness & Real-time correction \\
Identity (7) & $\uparrow W$ & Willingness & Deep preference shift \\
Temporal (8) & Optimal timing & All components & Timing leverage \\
\bottomrule
\end{tabular}
\end{center}

\textbf{Integration:} See HHH Axiom HHH-A3 for journey alignment and HHH-A5 for segment specificity.

\end{tcolorbox}

%------------------------------------------------------------------------------
% NATURAL TRAJECTORY MODEL (NTM) - NEW JANUARY 2026
%------------------------------------------------------------------------------
\subsection{Natural Trajectory Model: Behavior Dynamics Without Intervention}
\label{sec:bc-ntm}

\begin{tcolorbox}[colback=red!5!white, colframe=red!75!black,
    title=\textbf{Definition NTM-D1: Natural Trajectory (Do-Nothing Baseline)}]

\textbf{Statement:} The Natural Trajectory Model (NTM) describes the autonomous evolution of behavioral states when no intervention is applied ($\mathcal{I} = \emptyset$).

\begin{equation}
\vec{S}(t+\Delta t) = f_{NTM}(\vec{S}(t), \Psi(t), \Delta t, \mathbf{0})
\label{eq:ntm-evolution}
\end{equation}

where:
\begin{itemize}[nosep]
\item $\vec{S}(t) = (A(t), W(t), \theta(t), N(t), \varphi_{BCJ}(t))$ = State vector at time $t$
\item $\Psi(t)$ = Context at time $t$ (may drift autonomously)
\item $\Delta t$ = Time horizon
\item $\mathbf{0}$ = Null intervention (no action taken)
\end{itemize}

\textbf{Key Insight:} Understanding the natural trajectory is essential for:
\begin{enumerate}[nosep]
\item Establishing baseline for intervention effectiveness measurement
\item Identifying when intervention is necessary vs. when natural dynamics suffice
\item Predicting deterioration without action (urgency assessment)
\item Calculating opportunity cost of delayed intervention
\end{enumerate}

\end{tcolorbox}

\begin{tcolorbox}[colback=blue!5!white, colframe=blue!75!black,
    title=\textbf{Axiom NTM-A1: Four Time Horizons of Natural Dynamics}]

\textbf{Statement:} The natural trajectory exhibits distinct dynamics across four time horizons:

\vspace{0.5em}
\textbf{1. INSTANT ($\Delta t = 0$): Status Quo Persistence}
\begin{equation}
\vec{S}(t) = \vec{S}(t-\epsilon) \quad \text{(Status quo bias dominates)}
\end{equation}
\begin{itemize}[nosep]
\item Awareness: $A(t) = A_0$ (unchanged)
\item Willingness: $W(t) = WAX_0 - \theta_0$ (unchanged)
\item Context: $\Psi(t) = \Psi_0$ (stable)
\item $P(BC) = \sigma(WEC_0)$ (baseline probability)
\end{itemize}

\vspace{0.5em}
\textbf{2. SHORT-TERM ($\Delta t$ = weeks): Awareness Decay Begins}
\begin{equation}
A(t + \Delta t_{short}) = A_0 \cdot e^{-\lambda \Delta t} \quad \text{(AWX-A58)}
\end{equation}
\begin{itemize}[nosep]
\item Awareness: $A \downarrow$ by 15--35\% (decay rate $\lambda \approx 0.08$--$0.12$)
\item Willingness: $WAX \downarrow$ (since $A \downarrow$)
\item Context: $\Psi$ begins to drift (market, social norms)
\item Habits: Old behaviors reinforce (AWX-A60 habituation)
\item $P(BC) \downarrow$ by 15--25\%
\end{itemize}

\vspace{0.5em}
\textbf{3. MEDIUM-TERM ($\Delta t$ = months): Compounding Effects}
\begin{equation}
\theta(t + \Delta t_{medium}) = \theta_0 + \delta_{procrastination} \cdot \Delta t \quad \text{(Threshold drift)}
\end{equation}
\begin{itemize}[nosep]
\item Awareness: $A \approx 0.3 \cdot A_0$ (severely reduced)
\item Threshold: $\theta \uparrow$ (procrastination accumulates)
\item Context: $\Psi$-drift compounds (positive or negative)
\item Identity: Lock-in to current self-concept deepens
\item $P(BC) \downarrow$ by 40--60\% OR $\uparrow$ (if positive $\Psi$-drift)
\end{itemize}

\vspace{0.5em}
\textbf{4. LONG-TERM ($\Delta t$ = years): New Equilibrium}
\begin{equation}
\lim_{\Delta t \to \infty} \vec{S}(t + \Delta t) = \vec{S}_{equilibrium} \quad \text{(New steady state)}
\end{equation}
\begin{itemize}[nosep]
\item Awareness: $A_{eq} \ll A_0$ (minimal residual awareness)
\item Path Dependence: Early non-action constrains future options
\item Context: $\Psi$ may be fundamentally different
\item Identity: ``So bin ich halt'' crystallization
\item $P(BC) \to P_{eq}$ (convergence to new steady-state, potentially much lower)
\end{itemize}

\end{tcolorbox}

\begin{tcolorbox}[colback=green!5!white, colframe=green!75!black,
    title=\textbf{Theorem NTM-T1: Natural Trajectory Decay Rates by Component}]

\textbf{Statement:} Different state components decay at different rates without intervention:

\begin{center}
\renewcommand{\arraystretch}{1.3}
\begin{tabular}{@{}llll@{}}
\toprule
\textbf{Component} & \textbf{Decay Function} & \textbf{Half-Life} & \textbf{Empirical Source} \\
\midrule
Awareness $A(t)$ & $A_0 \cdot e^{-\lambda t}$ & 4--8 weeks & \citet{gabaix2019behavioral, sims2003implications} \\
Willingness $W(t)$ & $W_0 \cdot e^{-\mu t}$ & 2--4 months & \citet{loewenstein2009present, ariely2002procrastination} \\
Intent salience & $(1-\eta)^n$ per exposure & 3--6 exposures & \citet{wood2016psychology, bordalo2012salience} \\
Threshold $\theta(t)$ & $\theta_0 + \delta \cdot t$ & N/A (increases) & \citet{akerlof1991procrastination, odonoghue2001choice} \\
\bottomrule
\end{tabular}
\end{center}

\textbf{Implication:} Awareness decays fastest ($\lambda > \mu$), making AWARE and READY interventions time-critical. Willingness (WHO, HOW) is more persistent but harder to build.

\end{tcolorbox}

%------------------------------------------------------------------------------
% MATHEMATICAL DERIVATION OF DECAY FUNCTIONS
%------------------------------------------------------------------------------
\begin{tcolorbox}[colback=white, colframe=black!75!black,
    title=\textbf{Mathematical Derivation: NTM Decay Functions}]

\textbf{1. Exponential Decay (Awareness, Willingness)}

The exponential decay model follows from the assumption that the rate of decay is proportional to the current level:

\begin{equation}
\frac{dA}{dt} = -\lambda \cdot A(t)
\label{eq:awareness-ode}
\end{equation}

\textbf{Derivation:} Separating variables and integrating:
\begin{align}
\frac{dA}{A} &= -\lambda \, dt \\
\int_{A_0}^{A(t)} \frac{dA'}{A'} &= -\lambda \int_0^t dt' \\
\ln A(t) - \ln A_0 &= -\lambda t \\
A(t) &= A_0 \cdot e^{-\lambda t}
\end{align}

\textbf{Half-Life Relationship:} Setting $A(t_{1/2}) = \frac{A_0}{2}$:
\begin{equation}
t_{1/2} = \frac{\ln 2}{\lambda} \approx \frac{0.693}{\lambda}
\label{eq:half-life}
\end{equation}

\textbf{Parameter Estimation from Half-Life:}
\begin{center}
\renewcommand{\arraystretch}{1.3}
\begin{tabular}{@{}llll@{}}
\toprule
\textbf{Component} & \textbf{Half-Life} & \textbf{Decay Rate} & \textbf{Calculation} \\
\midrule
Awareness $A(t)$ & 4--8 weeks & $\lambda \approx 0.087$--$0.173$ /week & $\lambda = \ln(2)/t_{1/2}$ \\
Willingness $W(t)$ & 2--4 months & $\mu \approx 0.173$--$0.347$ /month & $\mu = \ln(2)/t_{1/2}$ \\
\bottomrule
\end{tabular}
\end{center}

\vspace{0.5em}
\textbf{2. Geometric Decay (Intent Salience)}

For discrete exposure-based decay, each exposure reduces salience by factor $(1-\eta)$:

\begin{equation}
S(n) = S_0 \cdot (1-\eta)^n
\label{eq:salience-decay}
\end{equation}

\textbf{Derivation:} After $n$ exposures without reinforcement:
\begin{align}
S(1) &= S_0 \cdot (1-\eta) \\
S(2) &= S(1) \cdot (1-\eta) = S_0 \cdot (1-\eta)^2 \\
&\vdots \\
S(n) &= S_0 \cdot (1-\eta)^n
\end{align}

\textbf{Exposure Half-Life:} Setting $S(n_{1/2}) = \frac{S_0}{2}$:
\begin{equation}
n_{1/2} = \frac{\ln 2}{\ln(1/(1-\eta))} = \frac{-\ln 2}{\ln(1-\eta)}
\label{eq:exposure-half-life}
\end{equation}

For $\eta = 0.15$--$0.25$: $n_{1/2} \approx 2.4$--$4.3$ exposures

\vspace{0.5em}
\textbf{3. Linear Threshold Drift (Procrastination)}

Procrastination accumulates linearly over time:

\begin{equation}
\frac{d\theta}{dt} = \delta \quad \Rightarrow \quad \theta(t) = \theta_0 + \delta \cdot t
\label{eq:threshold-drift}
\end{equation}

\textbf{Interpretation:} Each unit of time adds $\delta$ to the action threshold. Unlike decay, this has no equilibrium---without intervention, barriers grow indefinitely.

\textbf{Critical Time:} Time until threshold exceeds willingness permanently:
\begin{equation}
t_{critical} = \frac{W_0 - \theta_0}{\delta - \mu \cdot W_0} \quad \text{(if } \delta > \mu \cdot W_0 \text{)}
\end{equation}

\vspace{0.5em}
\textbf{4. Combined State Evolution}

The full NTM state vector evolves according to:
\begin{equation}
\vec{S}(t) = \begin{pmatrix} A_0 \cdot e^{-\lambda t} \\ W_0 \cdot e^{-\mu t} \\ \theta_0 + \delta t \\ N(t) \\ \varphi(t) \end{pmatrix}
\label{eq:state-evolution}
\end{equation}

\textbf{Behavioral Change Probability under NTM:}
\begin{equation}
P_{NTM}(t) = \sigma\left( A_0 e^{-\lambda t} \cdot (W_0 e^{-\mu t} - \theta_0 - \delta t) \cdot N(t) \right)
\label{eq:p-ntm}
\end{equation}

where $\sigma(\cdot)$ is the sigmoid function. Note: $P_{NTM}(t)$ decreases monotonically if $\delta > 0$ and no intervention occurs.

\end{tcolorbox}

\begin{tcolorbox}[colback=gray!5!white, colframe=gray!75!black,
    title=\textbf{Parameter Reference Table: Empirical Ranges}]

\begin{center}
\renewcommand{\arraystretch}{1.3}
\begin{tabular}{@{}lllll@{}}
\toprule
\textbf{Parameter} & \textbf{Symbol} & \textbf{Range} & \textbf{Unit} & \textbf{Source} \\
\midrule
Awareness decay rate & $\lambda$ & 0.087--0.173 & /week & Gabaix 2019, Sims 2003 \\
Willingness decay rate & $\mu$ & 0.173--0.347 & /month & Loewenstein 2009 \\
Salience attenuation & $\eta$ & 0.15--0.25 & /exposure & Bordalo 2012 \\
Threshold drift & $\delta$ & 0.01--0.05 & /month & O'Donoghue 2001 \\
Context drift (mean) & $\mu_\Psi$ & $-0.1$ to $+0.1$ & /month & Domain-specific \\
Context volatility & $\sigma_\Psi$ & 0.05--0.20 & /month & Domain-specific \\
\bottomrule
\end{tabular}
\end{center}

\textbf{Usage Notes:}
\begin{itemize}[nosep]
\item Parameters are \textit{population means}; individual variation is substantial
\item Domain-specific calibration recommended (see Appendix BBB for methods)
\item Context parameters ($\mu_\Psi$, $\sigma_\Psi$) require environmental analysis
\end{itemize}

\end{tcolorbox}

%------------------------------------------------------------------------------
% SCIENTIFIC FOUNDATIONS OF NTM
%------------------------------------------------------------------------------
\begin{tcolorbox}[colback=blue!5!white, colframe=blue!75!black,
    title=\textbf{Scientific Foundations: NTM in Historical Context}]

The NTM decay functions are not novel mathematical constructs but applications of well-established scientific models to behavioral dynamics. Each component has a direct historical precedent:

\vspace{0.5em}
\textbf{1. Awareness Decay $\leftarrow$ Ebbinghaus Forgetting Curve (1885)}

Hermann Ebbinghaus's seminal work established that memory retention follows exponential decay \citep{ebbinghaus1885memory}:
\begin{equation}
R(t) = e^{-t/S}
\end{equation}
where $R$ is retention and $S$ is memory strength. Our awareness function $A(t) = A_0 e^{-\lambda t}$ is mathematically identical, with $\lambda = 1/S$. The Ebbinghaus curve has been replicated across 140+ years of psychological research and remains the canonical model of forgetting.

\vspace{0.5em}
\textbf{2. Willingness Decay $\leftarrow$ Hyperbolic Discounting (Ainslie 1975, Laibson 1997)}

George Ainslie demonstrated that temporal discounting follows a hyperbolic rather than exponential form \citep{PAP-ainslie1975specious}:
\begin{equation}
D(t) = \frac{1}{1 + kt} \quad \text{(Hyperbolic)}
\end{equation}

David Laibson's quasi-hyperbolic $(\beta, \delta)$ model provides a tractable approximation \citep{laibson1997golden}:
\begin{equation}
U = u_0 + \beta \sum_{t=1}^{T} \delta^t u_t \quad \text{where } \beta < 1
\end{equation}

Our exponential willingness decay $W(t) = W_0 e^{-\mu t}$ captures the same phenomenon: motivation to act diminishes over time, with stronger discounting of distant outcomes.

\vspace{0.5em}
\textbf{3. Threshold Dynamics $\leftarrow$ Granovetter Threshold Models (1978)}

Mark Granovetter's threshold model of collective behavior established that individuals have heterogeneous action thresholds \citep{granovetter1978threshold}:
\begin{quote}
``An individual's threshold is the proportion of the group he would have to see join before he would do so.''
\end{quote}

Our dynamic threshold $\theta(t) = \theta_0 + \delta t$ extends this to individual-level temporal dynamics: without intervention, the psychological barrier to action increases through procrastination accumulation \citep{akerlof1991procrastination, odonoghue2001choice}.

\vspace{0.5em}
\textbf{4. Context Drift $\leftarrow$ Brownian Motion / Wiener Process (1923)}

The stochastic context evolution follows the standard form of geometric Brownian motion from physics and mathematical finance:
\begin{equation}
dX_t = \mu \, dt + \sigma \, dW_t
\end{equation}
where $W_t$ is a Wiener process. Our discrete formulation $\Psi(t+\Delta t) = \Psi(t) + \mu_\Psi \Delta t + \sigma_\Psi \epsilon(t)$ is the Euler-Maruyama discretization of this continuous process.

\vspace{0.5em}
\textbf{5. Salience Decay $\leftarrow$ Weber-Fechner Law \& Adaptation Level Theory}

The geometric decay of salience $(1-\eta)^n$ reflects established psychophysical principles: repeated exposure reduces perceived intensity (Weber-Fechner), and attention adapts to baseline levels (Helson's Adaptation Level Theory, 1964). Bordalo et al.'s salience theory \citep{PAP-bordalo2012salience} formalizes this for economic choice.

\end{tcolorbox}

\begin{tcolorbox}[colback=green!5!white, colframe=green!75!black,
    title=\textbf{Theorem NTM-T3: Scientific Validity of Decay Forms}]

\textbf{Statement:} The NTM decay functions are not ad-hoc but inherit validity from their scientific foundations.

\begin{center}
\renewcommand{\arraystretch}{1.3}
\begin{tabular}{@{}llll@{}}
\toprule
\textbf{NTM Component} & \textbf{Scientific Precedent} & \textbf{Year} & \textbf{Validation} \\
\midrule
Awareness decay & Ebbinghaus Forgetting Curve & 1885 & 140+ years replication \\
Willingness decay & Ainslie-Laibson $(\beta,\delta)$ & 1975/1997 & Nobel-adjacent (Thaler 2017) \\
Threshold drift & Granovetter Thresholds & 1978 & Sociology canon \\
Context drift & Wiener/Brownian Motion & 1923 & Physics foundation \\
Salience decay & Weber-Fechner / Bordalo & 1860/2012 & Psychophysics canon \\
\bottomrule
\end{tabular}
\end{center}

\textbf{Epistemic Status:} The NTM synthesizes established models into a unified framework for behavioral trajectory prediction. Individual components have Tier-1 (canonical) status; the synthesis is Tier-2 (theoretically derived, awaiting direct empirical validation of combined model).

\end{tcolorbox}

\begin{tcolorbox}[colback=yellow!5!white, colframe=yellow!75!black,
    title=\textbf{Axiom NTM-A2: Context Drift Without Intervention}]

\textbf{Statement:} Context $\Psi$ evolves autonomously according to external dynamics:

\begin{equation}
\Psi(t + \Delta t) = \Psi(t) + \underbrace{\mu_{\Psi} \cdot \Delta t}_{\text{Systematic drift}} + \underbrace{\sigma_{\Psi} \cdot \epsilon(t)}_{\text{Stochastic shock}}
\end{equation}

where:
\begin{itemize}[nosep]
\item $\mu_{\Psi}$ = Mean drift rate (can be positive or negative)
\item $\sigma_{\Psi}$ = Volatility of context shocks
\item $\epsilon(t) \sim N(0,1)$ = Random shock term
\end{itemize}

\textbf{Context Drift Categories:}
\begin{center}
\renewcommand{\arraystretch}{1.2}
\begin{tabular}{@{}lll@{}}
\toprule
\textbf{Drift Type} & \textbf{$\mu_{\Psi}$} & \textbf{Example} \\
\midrule
Positive drift & $> 0$ & Technology adoption, social norm shift \\
Neutral drift & $\approx 0$ & Stable market conditions \\
Negative drift & $< 0$ & Economic deterioration, trust erosion \\
High volatility & High $\sigma$ & Crisis, disruption, policy change \\
\bottomrule
\end{tabular}
\end{center}

\textbf{Implication:} Even without intervention, $P(BC)$ may increase (favorable drift) or decrease (unfavorable drift). The NTM must account for context uncertainty.

\end{tcolorbox}

%------------------------------------------------------------------------------
% EXTENDED CONTEXT DRIFT THEORY
%------------------------------------------------------------------------------
\begin{tcolorbox}[colback=orange!5!white, colframe=orange!75!black,
    title=\textbf{Definition NTM-D2: Multi-Dimensional Context Vector}]

\textbf{Statement:} Context $\Psi$ is not a scalar but a vector of distinct dimensions, each with its own drift dynamics:

\begin{equation}
\vec{\Psi}(t) = \begin{pmatrix} \Psi_{S}(t) \\ \Psi_{T}(t) \\ \Psi_{E}(t) \\ \Psi_{I}(t) \\ \Psi_{C}(t) \end{pmatrix} = \begin{pmatrix} \text{Social norms} \\ \text{Technology} \\ \text{Economic conditions} \\ \text{Institutions} \\ \text{Culture} \end{pmatrix}
\label{eq:psi-vector}
\end{equation}

Each dimension evolves according to its own stochastic process:
\begin{equation}
d\Psi_k = \mu_k \, dt + \sigma_k \, dW_k + \kappa_k \cdot \bar{B}(t) \, dt
\end{equation}

where $\bar{B}(t)$ is aggregate behavior (feedback term) and $\kappa_k$ is feedback strength.

\end{tcolorbox}

\begin{tcolorbox}[colback=cyan!5!white, colframe=cyan!75!black,
    title=\textbf{Axiom NTM-A3: Context Dimension Taxonomy}]

\textbf{Statement:} Each context dimension has distinct dynamics governed by established scientific theories:

\vspace{0.5em}
\begin{center}
\renewcommand{\arraystretch}{1.4}
\begin{tabular}{@{}p{1.8cm}p{3.5cm}p{2.5cm}p{3.5cm}@{}}
\toprule
\textbf{Dimension} & \textbf{Dynamics} & \textbf{Time Scale} & \textbf{Scientific Foundation} \\
\midrule
$\Psi_S$ (Social) & Norm evolution, peer effects & Weeks--Months & Bicchieri (2005), Cialdini (2004) \\
$\Psi_T$ (Tech) & S-curve adoption & Months--Years & Rogers (1962), Bass Model \\
$\Psi_E$ (Econ) & Market cycles, prices & Days--Months & Standard economics \\
$\Psi_I$ (Inst) & Rule/regulation change & Years--Decades & North (1990), Acemoglu \\
$\Psi_C$ (Culture) & Value/belief evolution & Decades--Generations & Boyd \& Richerson (1985) \\
\bottomrule
\end{tabular}
\end{center}

\textbf{Hierarchy of Time Scales:}
\begin{equation}
\tau_E < \tau_S < \tau_T < \tau_I < \tau_C
\end{equation}

Economic context changes fastest; cultural context changes slowest. Intervention design must match the relevant time scale.

\end{tcolorbox}

\begin{tcolorbox}[colback=red!5!white, colframe=red!75!black,
    title=\textbf{Axiom NTM-A4: Behavior-Context Feedback Loop}]

\textbf{Statement:} Aggregate behavior feeds back into context, creating potential for self-reinforcing dynamics:

\begin{equation}
\frac{d\Psi_S}{dt} = \underbrace{\mu_S}_{\text{Exogenous drift}} + \underbrace{\kappa \cdot (\bar{B}(t) - \bar{B}^*)}_{\text{Endogenous feedback}} + \sigma_S \cdot \epsilon(t)
\label{eq:feedback-loop}
\end{equation}

where:
\begin{itemize}[nosep]
\item $\bar{B}(t) = \frac{1}{N}\sum_i B_i(t)$ = Population mean behavior
\item $\bar{B}^*$ = Reference behavior level (status quo)
\item $\kappa > 0$ = Feedback strength (social multiplier)
\end{itemize}

\textbf{Mechanism:} When more people adopt a behavior ($\bar{B} > \bar{B}^*$), social norms shift favorably ($d\Psi_S > 0$), which further increases adoption probability. This creates positive feedback.

\textbf{Evidence:}
\begin{itemize}[nosep]
\item Social proof effects \citep{cialdini2004social}
\item Descriptive norm influence \citep{bicchieri2005grammar}
\item Network externalities \citep{PAP-rogers2003diffusion}
\end{itemize}

\end{tcolorbox}

\begin{tcolorbox}[colback=yellow!5!white, colframe=yellow!75!black,
    title=\textbf{Theorem NTM-T4: Tipping Point Dynamics (Granovetter-Schelling)}]

\textbf{Statement:} Context drift exhibits threshold effects---below critical mass, change is slow; above it, change accelerates dramatically.

\begin{equation}
\frac{d\Psi_S}{dt} = \begin{cases}
\mu_{slow} & \text{if } \bar{B}(t) < \theta_{tip} \\[0.5em]
\mu_{fast} & \text{if } \bar{B}(t) \geq \theta_{tip}
\end{cases}
\quad \text{where } \mu_{fast} \gg \mu_{slow}
\label{eq:tipping}
\end{equation}

\textbf{Derivation from Granovetter (1978):}

Let $F(\theta)$ be the cumulative distribution of individual thresholds in the population. An individual adopts when the proportion of adopters exceeds their personal threshold. The equilibrium condition:
\begin{equation}
\bar{B}^* = F(\bar{B}^*)
\end{equation}

\textbf{Tipping occurs when:}
\begin{equation}
\frac{dF}{d\bar{B}}\bigg|_{\bar{B}=\theta_{tip}} > 1
\end{equation}

At this point, a small increase in adoption triggers a cascade.

\textbf{Empirical Tipping Points:}
\begin{center}
\renewcommand{\arraystretch}{1.2}
\begin{tabular}{@{}lll@{}}
\toprule
\textbf{Domain} & \textbf{$\theta_{tip}$} & \textbf{Source} \\
\midrule
Social norm change & 25--35\% & Centola et al. (2018) \\
Technology adoption & 10--20\% (innovators + early adopters) & Rogers (1962) \\
Opinion cascades & 10--15\% committed minority & Xie et al. (2011) \\
\bottomrule
\end{tabular}
\end{center}

\textbf{Implication for NTM:} Without intervention, natural drift may never reach $\theta_{tip}$. Strategic intervention can ``push'' the system past the tipping point, after which favorable context drift becomes self-sustaining.

\end{tcolorbox}

\begin{tcolorbox}[colback=green!5!white, colframe=green!75!black,
    title=\textbf{Theorem NTM-T5: Context Drift Scientific Foundations}]

\textbf{Statement:} Each context dimension's drift dynamics are grounded in canonical scientific theories:

\vspace{0.5em}
\textbf{1. Social Norm Drift ($\Psi_S$) $\leftarrow$ Bicchieri \& Cialdini}

Social norms evolve through changes in empirical expectations (what others do) and normative expectations (what others approve) \citep{bicchieri2005grammar}:
\begin{equation}
\Psi_S(t+1) = f(\text{Empirical}_t, \text{Normative}_t, \text{Sanctions}_t)
\end{equation}

\textbf{2. Technology Drift ($\Psi_T$) $\leftarrow$ Rogers Diffusion Model (1962)}

Technology adoption follows an S-curve driven by innovators → early adopters → majority → laggards \citep{PAP-rogers2003diffusion}:
\begin{equation}
\frac{dN}{dt} = p(M - N) + q\frac{N}{M}(M - N) \quad \text{(Bass Model)}
\end{equation}

\textbf{3. Institutional Drift ($\Psi_I$) $\leftarrow$ North (1990), Acemoglu}

Institutions change through:
\begin{itemize}[nosep]
\item Formal rules (legislation, regulation)
\item Informal constraints (conventions, norms)
\item Enforcement mechanisms
\end{itemize}
Institutional change is typically slow but can exhibit punctuated equilibrium \citep{PAP-north1990institutions}.

\textbf{4. Cultural Drift ($\Psi_C$) $\leftarrow$ Boyd \& Richerson (1985)}

Cultural evolution operates through:
\begin{equation}
\Psi_C(t+1) = (1-\mu)\Psi_C(t) + \mu \cdot \text{Transmission}(\text{Parents}, \text{Peers}, \text{Institutions})
\end{equation}
where $\mu$ is the cultural mutation/innovation rate \citep{boyd1985culture}.

\textbf{5. Path Dependence $\leftarrow$ Arthur (1994)}

All context dimensions exhibit path dependence---history matters \citep{arthur1994increasing}:
\begin{equation}
\Psi(t) = f(\Psi(0), \{\epsilon_s\}_{s=0}^{t})
\end{equation}
Early random events can lock in long-term trajectories.

\end{tcolorbox}

\begin{tcolorbox}[colback=gray!5!white, colframe=gray!75!black,
    title=\textbf{Parameter Table: Context Drift by Dimension}]

\begin{center}
\renewcommand{\arraystretch}{1.3}
\begin{tabular}{@{}llllll@{}}
\toprule
\textbf{Dimension} & \textbf{$\mu_k$} & \textbf{$\sigma_k$} & \textbf{$\kappa_k$} & \textbf{$\tau_k$} & \textbf{Source} \\
\midrule
$\Psi_S$ (Social) & $\pm 0.02$/mo & 0.05/mo & 0.3--0.5 & Weeks & Cialdini \\
$\Psi_T$ (Tech) & $+0.01$/mo & 0.03/mo & 0.2--0.4 & Months & Rogers \\
$\Psi_E$ (Econ) & $\pm 0.05$/mo & 0.15/mo & 0.1 & Days & Market data \\
$\Psi_I$ (Inst) & $\pm 0.005$/yr & 0.02/yr & 0.05 & Years & North \\
$\Psi_C$ (Culture) & $\pm 0.001$/yr & 0.005/yr & 0.02 & Decades & Boyd \\
\bottomrule
\end{tabular}
\end{center}

\textbf{Usage:} These are order-of-magnitude estimates. Domain-specific calibration required via Appendix BBB methods.

\end{tcolorbox}

\begin{tcolorbox}[colback=purple!5!white, colframe=purple!75!black,
    title=\textbf{Theorem NTM-T2: Intervention Necessity Criterion}]

\textbf{Statement:} Intervention is necessary if and only if the natural trajectory leads to an unacceptable state:

\begin{equation}
\text{Intervention Required} \Leftrightarrow P_{NTM}(t + \Delta t) < P_{target} - \epsilon
\end{equation}

where:
\begin{itemize}[nosep]
\item $P_{NTM}(t + \Delta t)$ = Probability under natural trajectory
\item $P_{target}$ = Desired behavioral change probability
\item $\epsilon$ = Acceptable deviation margin
\end{itemize}

\textbf{Decision Rule:}
\begin{center}
\renewcommand{\arraystretch}{1.2}
\begin{tabular}{@{}ll@{}}
\toprule
\textbf{Condition} & \textbf{Action} \\
\midrule
$P_{NTM} \geq P_{target}$ & No intervention needed (natural dynamics sufficient) \\
$P_{NTM} < P_{target}$, gap small & Light intervention (AWARE, READY, WHEN) \\
$P_{NTM} \ll P_{target}$ & Heavy intervention (WHO, WHAT$_S$, WHAT$_F$, HOW) \\
$P_{NTM} \to 0$ rapidly & Urgent intervention with portfolio approach \\
\bottomrule
\end{tabular}
\end{center}

\end{tcolorbox}

\begin{tcolorbox}[colback=gray!5!white, colframe=gray!75!black,
    title=\textbf{Scientific Evidence for NTM Components}]

\textbf{The Natural Trajectory Model synthesizes 50+ years of behavioral economics research:}

\vspace{0.5em}
\textbf{1. Status Quo Persistence (NTM-D1):}
\begin{itemize}[nosep]
\item Status quo bias and endowment effect \citep{kahneman1991anomalies, samuelson1988status}
\item Loss aversion amplifies inertia: $\lambda \approx 2$ \citep{kahneman1979prospect, tversky1991loss}
\end{itemize}

\vspace{0.5em}
\textbf{2. Awareness Decay (NTM-A1, Short-term):}
\begin{itemize}[nosep]
\item Hyperbolic discounting and present bias \citep{laibson1997golden, frederick2002time, loewenstein2009present}
\item Behavioral inattention and limited attention \citep{gabaix2019behavioral, sims2003implications}
\item Salience decay without reinforcement \citep{PAP-bordalo2012salience}
\end{itemize}

\vspace{0.5em}
\textbf{3. Threshold Drift and Procrastination (NTM-A1, Medium-term):}
\begin{itemize}[nosep]
\item Naive present bias leads to procrastination \citep{odonoghue2001choice, akerlof1991procrastination}
\item Contract design exploiting time inconsistency \citep{PAP-dellavigna2004contract}
\item Intent-action gap widens over time \citep{ajzen1991theory, sheeran2016intention}
\end{itemize}

\vspace{0.5em}
\textbf{4. Habit Lock-In and Identity Crystallization (NTM-A1, Long-term):}
\begin{itemize}[nosep]
\item Habit formation and automaticity \citep{wood2016psychology, verplanken2006habit}
\item Identity-based preferences become stable \citep{PAP-akerlof2000economics}
\item Mental accounting budgets crystallize \citep{thaler1985mental, thaler1999mental}
\end{itemize}

\vspace{0.5em}
\textbf{5. Context Drift (NTM-A2):}
\begin{itemize}[nosep]
\item Social norm evolution \citep{bicchieri2005grammar, cialdini2004social}
\item Market and technology dynamics \citep{PAP-rogers2003diffusion}
\item Policy and regulatory changes \citep{thaler2008nudge, sunstein2014nudging}
\end{itemize}

\textbf{Meta-Evidence:} NTM decay rates validated by meta-analyses showing intervention effect sizes decay 30--50\% from pilot to scale \citep{hummel2019effective, dellavigna2022rcts, soman2020successfully}.

\end{tcolorbox}


%===============================================================================
% UNIFIED NTM-CONTEXT MODEL (UNTCM) - NOW IN DEDICATED APPENDIX
%===============================================================================

\subsection{Unified NTM-Context Model (UNTCM)}
\label{subsec:untcm}

\begin{tcolorbox}[colback=orange!5!white, colframe=orange!75!black,
    title=\textbf{Cross-Reference: UNTCM moved to Appendix UN}]

\textbf{The complete Unified NTM-Context Model (UNTCM) is now documented in:}

\begin{center}
\Large\textbf{Appendix UN: FORMAL-UNTCM}
\end{center}

\textbf{UNTCM Content in Appendix UN:}
\begin{itemize}[nosep]
    \item \textbf{Core Theory:} A5 (Context-Modulated Decay), T6 (Coupled ODEs), T7 (Phase Diagram), T8 (Timing Rule)
    \item \textbf{Methodology:} M1 (Computational), M2 (Calibration), M3 (Stochastic), M4 (Optimal Control)
    \item \textbf{Integration:} I1 (BCJ), I2 (T1-T8), I3 (FEPSDE), I4 (Segments)
    \item \textbf{Empirical:} E1 (Regime Detection), E2 (Tipping Prediction), E3 (Timing Estimation)
    \item \textbf{Extensions:} X1 (Network), X2 (Hysteresis), X3 (Adaptive $\theta$), X4 (Markov-Switching)
    \item \textbf{Evidence Base:} 22+ PRO papers, 4 constructive CONTRA papers
\end{itemize}

\textbf{Why the Move?}
\begin{itemize}[nosep]
    \item UNTCM grew to 25 elements (A5, T6-T8, M1-M4, I1-I4, E1-E3, X1-X4) with extensive evidence base
    \item Dedicated appendix provides better navigation and visibility
    \item PBB remains focused on probability axioms (P1-P10) and foundational NTM (A1-A4, T1-T5)
\end{itemize}

\textbf{Key Connection:} PBB provides the foundational NTM axioms (A1-A4) and theorems (T1-T5). Appendix UN extends this to the \emph{unified} model that couples NTM dynamics with context evolution.

\vspace{0.5em}
\textit{See Appendix UN (FORMAL-UNTCM) for the complete formalization.}

\end{tcolorbox}

%===============================================================================
% OPEN ISSUES
%===============================================================================

\section{Open Issues and Research Frontiers}

\begin{enumerate}

\item \textbf{Heterogeneous Learning Rates:} Different individuals learn at different rates from experience. Current framework uses homogeneous $\eta(\text{Learning})$; this should be segment-dependent.

\item \textbf{Feedback Loop Dynamics:} When behavior changes, does this change segment membership? When context changes, how quickly do probabilities re-equilibrate? Dynamic modeling needed.

\item \textbf{Cascade Effects:} When one person transitions, does this increase probability for peers through social norm effects? Cascade modeling not yet incorporated.

\item \textbf{Irreversibility:} Some transitions are hard to reverse (e.g., solar adoption is durable; smoking cessation often reverses). How to incorporate asymmetric reversibility?

\item \textbf{Higher-Order Interactions:} Current framework assumes additive (multiplicative) combination of terms. Should we model higher-order interactions between willingness and utility, or between segment and context?

\end{enumerate}

%===============================================================================
% SUMMARY
%===============================================================================

\section{Summary: Why These 10 Axioms Matter}

The 10 axioms formalize the **probabilistic nature of behavior change**:

\begin{itemize}[nosep]

\item \textbf{Axiom 1:} Ensures all BCM components integrate into transition probabilities

\item \textbf{Axiom 2:} Captures utility's positive effect

\item \textbf{Axiom 3:} Prevents false predictions when willingness is blocked (critical practical insight)

\item \textbf{Axiom 4:} Shows how awareness activates in context

\item \textbf{Axiom 5:} Formalizes context dependency

\item \textbf{Axiom 6:} Prevents one-size-fits-all models; ensures segment differentiation

\item \textbf{Axiom 7:} Captures learning and habituation dynamics

\item \textbf{Axiom 8:} Incorporates social effects

\item \textbf{Axiom 9:} Prevents overconfidence in simple incentive models; shows risk/uncertainty matter

\item \textbf{Axiom 10:} Ensures models are testable, not merely theoretical

\end{itemize}

Together, they enable:
\begin{itemize}[nosep]
    \item Transparent, mathematically rigorous modeling
    \item Empirically testable predictions
    \item Segment-specific intervention design
    \item Realistic probability estimation that respects behavioral complexity
\end{itemize}

%===============================================================================
% GLOSSARY REFERENCE
%===============================================================================

\section{Glossary Reference}

For definitions of key terms used in this appendix, see **Appendix GLS (REF-GLOSSARY)** which provides formal definitions for:
\begin{itemize}[nosep]
    \item \textbf{Probability:} Mathematical definition and interpretations (frequentist vs. subjective)
    \item \textbf{Utility:} Net utility combining all FEPSDE dimensions
    \item \textbf{Willingness:} Dispositional readiness to engage with behavior
    \item \textbf{Threshold:} Individual-specific barrier to action
    \item \textbf{Segment:} Behavioral change segment classification (Innovator, Pragmatist, Risk-Averse, Laggard)
    \item \textbf{Journey Stage:} Phase in behavioral change journey (Unaware → Aware → Considering → Intending → Acting → Maintaining)
    \item \textbf{Context (Ψ):} Situational and structural factors affecting behavior
\end{itemize}

%===============================================================================
% REFERENCES
%===============================================================================

\section{References}

This appendix integrates research from the master bibliography maintained in the BEATRIX repository. Key references include:

\cite{bcm_master}

All theoretical claims in Axioms P1-P10 are grounded in peer-reviewed literature including:
\begin{itemize}[nosep]
    \item Fehr & Gächter (50+ papers on behavioral economics and preferences)
    \item Kahneman & Tversky (prospect theory, decision-making under uncertainty)
    \item Thaler & Sunstein (behavioral economics, nudges, choice architecture)
    \item Ryan & Deci (self-determination theory, intrinsic motivation)
    \item Ajzen (theory of planned behavior, intentions-behavior gap)
    \item Rogers (diffusion of innovations, adoption curves)
    \item Cialdini (social influence principles)
\end{itemize}

For full citations and database of 1000+ papers supporting EBF framework, see Master Bibliography reference.

%===============================================================================

\end{document}
